% !TeX program = pdflatex
%
% Map2Check v8 -- Modos, propriedades, sub-propriedades e cobertura de CWEs
%
% Documento de referencia levantado diretamente do codigo-fonte da branch
% wasm-verification-upgrade. Cada tabela cita o arquivo e a linha de onde o
% dado foi extraido, para que qualquer divergencia futura seja rastreavel.
%
% Compilar com pdfLaTeX (duas passadas). Sem dependencias fora do TeX Live
% scheme-small.
%
\documentclass[11pt,a4paper]{article}

\usepackage[T1]{fontenc}
\usepackage[utf8]{inputenc}
\usepackage[brazil]{babel}
\usepackage{lmodern}
\usepackage{microtype}

\setlength{\emergencystretch}{3.5em}

\usepackage[margin=2.5cm]{geometry}
\usepackage{booktabs}
\usepackage{tabularx}
\usepackage{longtable}
\setlength{\LTcapwidth}{\textwidth}
\usepackage{array}
\usepackage{xcolor}
\usepackage{enumitem}
\usepackage{float}
\usepackage{fancyhdr}
\usepackage[hidelinks]{hyperref}
\urlstyle{tt}
\def\UrlBreaks{\do\/\do\.\do\-\do\_\do\{\do\}}

%---------------------------------------------------------------- macros
\newcommand{\enquote}[1]{``#1''}
\newcommand{\code}[1]{\texttt{#1}}
\newcommand{\flag}[1]{\texttt{-\/-#1}}
\newcommand{\src}[1]{{\small\texttt{#1}}}
\newcommand{\yes}{\textcolor{green!45!black}{\textbf{sim}}}
\newcommand{\no}{\textcolor{red!65!black}{\textbf{não}}}
\newcommand{\parcial}{\textcolor{orange!75!black}{\textbf{parcial}}}

\newcolumntype{L}[1]{>{\raggedright\arraybackslash}p{#1}}
\newcolumntype{C}[1]{>{\centering\arraybackslash}p{#1}}
\newcolumntype{Y}{>{\raggedright\arraybackslash}X}

\pagestyle{fancy}
\fancyhf{}
\fancyhead[L]{\small Map2Check --- Propriedades e CWEs}
\fancyhead[R]{\small\thepage}
\renewcommand{\headrulewidth}{0.4pt}

\title{\textbf{Map2Check}\\[2mm]
\large Modos de verificação, propriedades, sub-propriedades\\ e cobertura de CWEs}
\author{Levantamento a partir do código-fonte\\
\small branch \code{wasm-verification-upgrade} --- Map2Check v8.0.0}
\date{6 de setembro de 2026}

\begin{document}
\maketitle

\begin{abstract}
\noindent
Este documento inventaria \emph{tudo} que o Map2Check consegue verificar hoje:
os seis modos de verificação implementados, as especificações SV-COMP que cada
um responde, os dez veredictos que o runtime é capaz de escrever, e o
mapeamento CWE${\rightarrow}$modo tal como ele existe nos dois harnesses de
avaliação do repositório (CASTLE e Juliet). Nada aqui é aspiracional: cada
entrada cita o arquivo e a linha que a sustenta, e as limitações conhecidas
estão na Seção~\ref{sec:limitacoes} em vez de omitidas.
\end{abstract}

\tableofcontents
\bigskip

%================================================================
\section{Escopo e fontes}

O levantamento foi feito sobre quatro camadas do código, porque a resposta
completa não está em nenhuma delas isoladamente:

\begin{enumerate}[leftmargin=*,itemsep=2pt]
  \item \textbf{CLI} --- \src{modules/frontend/map2check.cpp} define as flags e
        traduz cada uma em um modo.
  \item \textbf{Orquestração} --- \src{modules/frontend/caller.hpp:21} declara
        \code{enum class Map2CheckMode}; \src{caller.cpp:389} decide quais
        \emph{passes} LLVM cada modo carrega.
  \item \textbf{Instrumentação} --- \src{modules/backend/pass/*.cpp} são os
        passes que injetam as chamadas de verificação no IR.
  \item \textbf{Runtime} --- \src{modules/backend/library/lib/AnalysisMode*.c}
        implementam a checagem em tempo de execução e escrevem o veredicto;
        \src{header/Map2CheckTypes.h:68} enumera os veredictos possíveis.
\end{enumerate}

A cobertura de CWEs vem dos dois \emph{runners} de avaliação versionados no
repositório, que são a única declaração executável desse mapeamento:
\src{tests/castle/run\_castle\_evaluation.sh} e
\src{tests/juliet/run\_juliet\_evaluation.sh}. O contrato entre eles é
verificado por \src{tests/integration/test\_cwe\_mode\_mapping.sh}, que impede
que qualquer CWE do benchmark caia num modo não declarado.

%================================================================
\section{Modos de verificação}

São seis modos, dos quais cinco checam uma propriedade e um --- Cover-Branches
--- não checa propriedade alguma, existindo apenas para produzir suítes de
teste.

\begin{table}[H]
\centering
\small
\begin{tabularx}{\textwidth}{@{}l L{3.4cm} L{3.75cm} Y@{}}
\toprule
\textbf{\#} & \textbf{Modo} & \textbf{Flag} & \textbf{O que checa} \\
\midrule
1 & \code{MEMTRACK\_MODE} \newline \emph{(padrão)} & \flag{memtrack} &
Segurança de memória: \emph{free} inválido, dereferência inválida e vazamento
de memória alocada. \\
\addlinespace
2 & \code{MEMCLEANUP\_MODE} & \flag{memcleanup-property} &
Toda memória alocada foi liberada até o fim do \code{main}. \\
\addlinespace
3 & \code{OVERFLOW\_MODE} & \flag{check-overflow} &
Overflow de inteiro \emph{com sinal} e divisão/resto por zero. \\
\addlinespace
4 & \code{REACHABILITY\_MODE} & \flag{target-function} \newline
\flag{target-function-name} & Se uma função-alvo nomeada é alcançável. \\
\addlinespace
5 & \code{ASSERT\_MODE} & \flag{check-asserts} &
Se algum \code{\_\_VERIFIER\_assert} ou \code{\_\_assert\_fail} pode falhar. \\
\addlinespace
6 & \code{COVER\_BRANCHES\_MODE} & \flag{cover-branches} &
\emph{Nada.} Explora caminhos e registra suas entradas (Test-Comp
Cover-Branches). \\
\bottomrule
\end{tabularx}
\caption{Os seis modos. Fonte: \src{caller.hpp:21--37} e
\src{map2check.cpp:781--878}.}
\end{table}

\subsection{Especificação SV-COMP correspondente}

É preciso separar duas coisas que o código trata em lugares diferentes, e que
não coincidem.

\paragraph{Na \emph{witness}.} A especificação é escolhida pela
\textbf{propriedade violada}, não pelo modo --- \src{witness/graph.cpp:208--237}
faz um \code{switch} sobre \code{PropertyViolated}.

\begin{table}[H]
\centering
\small
\begin{tabularx}{\textwidth}{@{}L{3.6cm} Y@{}}
\toprule
\textbf{Veredicto} & \textbf{Especificação emitida} \\
\midrule
\code{FALSE-FREE}     & \code{CHECK( init(main()), LTL(G valid-free) )} \\
\code{FALSE-DEREF}    & \code{CHECK( init(main()), LTL(G valid-deref) )} \\
\code{FALSE-MEMTRACK} & \code{CHECK( init(main()), LTL(G valid-memtrack) )} \\
\code{FALSE-OVERFLOW} & \code{CHECK( init(main()), LTL(G ! overflow) )} \\
\code{TARGET-REACHED} & especificação de alvo (\code{SpecificationType::TARGET}) \\
\addlinespace
\code{NONE} (\emph{TRUE}) & conforme \code{spectTrue}: alvo, \code{! overflow},
ou --- por padrão --- as três linhas de \code{valid-free} $+$
\code{valid-deref} $+$ \code{valid-memtrack} (\code{MEMSAFETY}). \\
\bottomrule
\end{tabularx}
\caption{Fonte: \src{witness/witness.hpp:64--101} (enum
\code{SpecificationType}) e \src{witness/graph.cpp:198--258}.
\textbf{Não existe especificação de \code{memcleanup}}: o enum
\code{SpecificationType} tem seis valores e nenhum deles é
\code{valid-memcleanup}. Ver a Seção~\ref{sec:limitacoes}, item~3.}
\end{table}

\paragraph{Nos metadados do \emph{test-suite}.} Aqui a especificação vem do
arquivo passado em \flag{property-file}, copiado \emph{verbatim}. As duas
strings abaixo são apenas o \emph{fallback} para quando esse arquivo não é
passado ou não pode ser lido (\src{map2check.cpp:104--128}):

\begin{table}[H]
\centering
\small
\begin{tabularx}{\textwidth}{@{}L{3.6cm} Y@{}}
\toprule
\textbf{Modo} & \textbf{\emph{Fallback} de \code{<specification>}} \\
\midrule
\code{REACHABILITY} &
\code{COVER( init(main()), FQL(COVER EDGES(@CALL(reach\_error))) )} \\
\addlinespace
demais modos &
\code{COVER( init(main()), FQL(COVER EDGES(@DECISIONEDGE)) )} \\
\bottomrule
\end{tabularx}
\end{table}

\subsection{Pipeline de passes por modo}

Cada modo monta uma sequência distinta de \emph{passes} do New Pass Manager. O
\code{nondet-pass} e o \code{map2check-library} são carregados sempre.

\begin{table}[H]
\centering
\small
\begin{tabularx}{\textwidth}{@{}L{3.4cm} L{4.3cm} Y@{}}
\toprule
\textbf{Modo} & \textbf{Pass adicional} & \textbf{Runtime} \\
\midrule
\code{MEMTRACK}       & \code{MemoryTrackPass}  & \code{AnalysisModeMemtrack.c} \\
\code{MEMCLEANUP}     & \code{MemoryTrackPass}  & \code{AnalysisModeMemcleanup.c} \\
\code{OVERFLOW}       & \code{OverflowPass}     & \code{AnalysisModeOverflow.c} \\
\code{REACHABILITY}   & \code{TargetPass}       & \code{AnalysisModeNone.c} \\
\code{ASSERT}         & \code{AssertPass}       & \code{AnalysisModeAssert.c} \\
\code{COVER\_BRANCHES}& \emph{nenhum}           & \code{AnalysisModeNone.c} \\
\bottomrule
\end{tabularx}
\caption{Passes: \src{caller.cpp:389--445}; \emph{runtime} linkado:
\src{caller.cpp:499--531}. \code{REACHABILITY} não precisa de um
\code{AnalysisMode} próprio porque \code{TARGET-REACHED} é escrito pela
biblioteca comum (\src{Map2CheckFunctions.c:62}). Cover-Branches é
deliberadamente o pipeline mais enxuto: menos instrumentação significa mais
caminhos explorados pelo KLEE no mesmo orçamento.}
\end{table}

%================================================================
\section{Veredictos e sub-propriedades}

Um modo não produz um único veredicto. O runtime escreve em
\code{map2check\_property} um dentre nove valores --- o \emph{enum} tem dez,
mas \code{UNKNOWN} nunca chega lá (ver abaixo) --- oito dos quais são violações
concretas. Esta é a granularidade real da ferramenta --- e é o nível certo para
uma comparação contra outra ferramenta, já que \code{FALSE-FREE},
\code{FALSE-DEREF} e \code{FALSE-MEMTRACK} compartilham a mesma flag mas são
defeitos distintos.

\begin{table}[H]
\centering
\small
\begin{tabularx}{\textwidth}{@{}c L{3.4cm} L{3.3cm} Y@{}}
\toprule
\textbf{Cód.} & \textbf{Veredicto} & \textbf{Emitido por} & \textbf{Significado} \\
\midrule
0 & \code{NONE}             & todos                & Nenhuma violação encontrada (\emph{TRUE}). \\
1 & \code{UNKNOWN} \newline {\footnotesize(nunca gravado)} & --- & Indeciso. \textbf{As duas chamadas que o escreveriam estão comentadas} (\src{AnalysisModeMemtrack.c:308}, \src{AnalysisModeMemcleanup.c:298}); ver a Seção~\ref{sec:limitacoes}, item~11. \\
2 & \code{FALSE-FREE}       & Memtrack, Memcleanup & \emph{free} de ponteiro inválido, \emph{double free}, \emph{free} de memória não alocada. Também o \code{free} de WASM sobre \emph{offset} não alocado. \\
3 & \code{FALSE-DEREF}      & Memtrack             & Dereferência de ponteiro inválido (nulo, liberado, fora da alocação). \\
4 & \code{FALSE-MEMTRACK}   & Memtrack             & Memória alocada e não liberada (\emph{leak}). \\
5 & \code{TARGET-REACHED}   & TargetPass           & A função-alvo é alcançável. \\
6 & \code{FALSE-OVERFLOW} \newline {\footnotesize(gravado como \code{OVERFLOW})} & Overflow; \newline Memtrack (WASM) & Overflow em aritmética de inteiro com sinal. \textbf{No modo WASM é também o veredicto de acesso fora dos limites} da \emph{linear memory} (\src{AnalysisModeMemtrack.c:418--424}). \\
7 & \code{ASSERT}           & AssertPass           & Assertiva do usuário violada. \\
8 & \code{FALSE-MEMCLEANUP} & Memcleanup           & Memória ainda viva no retorno do \code{main}. \\
9 & \code{FALSE-DIVBYZERO}  & Overflow             & Divisão ou resto por zero. \\
\bottomrule
\end{tabularx}
\caption{Fonte: \src{library/header/Map2CheckTypes.h:68--79} (enum
\code{ViolatedProperty}) e \src{frontend/utils/tools.hpp:105--116}
(\code{PropertyViolated}). Os pontos de emissão estão em
\src{AnalysisModeMemtrack.c:141,174,299,356,413,421},
\src{AnalysisModeMemcleanup.c:142,170,346},
\src{AnalysisModeOverflow.c:193} e \src{AnalysisModeAssert.c:20}.}
\end{table}

\paragraph{Observação importante.} \code{FALSE-DIVBYZERO} \emph{não tem flag
própria}: divisão por zero é detectada como efeito colateral de
\flag{check-overflow}, porque a instrumentação de \code{SDiv}/\code{SRem} é a
mesma que checa o overflow de \code{INT\_MIN / -1}. Qualquer comparação com
outra ferramenta precisa acionar \flag{check-overflow} para cobrir CWE-369.

%================================================================
\section{Operações instrumentadas em \texttt{OVERFLOW\_MODE}}

\begin{table}[H]
\centering
\small
\begin{tabularx}{\textwidth}{@{}L{3.3cm} C{2.9cm} Y@{}}
\toprule
\textbf{Instrução LLVM} & \textbf{Instrumentada} & \textbf{Nota} \\
\midrule
\code{Add}, \code{Sub}, \code{Mul} & \yes &
Apenas a variante \emph{com sinal}. O critério é a \emph{flag} \code{nsw} da
instrução: \src{isUnitAssignment = !binOp->hasNoSignedWrap()}, então só a
aritmética que o clang marcou como \code{nsw} é instrumentada. As versões
\code{uint} existem no código porém estão comentadas. \\
\code{SDiv} & \yes & Gera \code{FALSE-OVERFLOW} (\code{INT\_MIN / -1}) ou \code{FALSE-DIVBYZERO}. \\
\code{SRem} & \yes & Idem. \\
\addlinespace
\code{UDiv}, \code{URem} & \no & Aritmética sem sinal não faz parte da propriedade \code{no-overflow}. \\
\code{Shl}, \code{LShr}, \code{AShr} & \no &
O \code{case} de \code{Shl} existe no \code{switch} mas tem \textbf{corpo
vazio} --- deslocamento à esquerda com perda de bits \emph{não} é detectado. \\
\code{And}, \code{Or}, \code{Xor} & \no & \emph{Case} vazio. \\
\code{FAdd}, \code{FSub}, \code{FMul}, \code{FDiv}, \code{FRem} & \no &
\emph{No-op} deliberado: o SV-COMP define \code{no-overflow} apenas sobre tipos
inteiros. \\
\bottomrule
\end{tabularx}
\caption{Fonte: \src{pass/OverflowPass.cpp:325--424}. São \textbf{cinco} as
instruções efetivamente instrumentadas. A propriedade é avaliada
sobre \emph{o tipo resultante da operação}, conforme a definição do SV-COMP ---
ver o comentário em \src{OverflowPass.cpp:418} e
\src{pass/OperationsFunctions.hpp:63}.}
\end{table}

%================================================================
\section{Cobertura de CWEs}

O repositório contém dois mapeamentos CWE${\rightarrow}$modo, ambos executáveis
e ambos versionados. Eles não são idênticos porque os benchmarks não são: o
CASTLE-C250 cobre CWEs de segurança mais amplas, e o Juliet cobre famílias de
memória e aritmética com muitas variantes por família.

\subsection{Mapeamento consolidado}

\begin{longtable}{@{\extracolsep{\fill}}l L{4.5cm} L{3.9cm} C{1.7cm} C{1.5cm}@{}}
\caption{União dos dois mapeamentos. \textbf{15 CWEs} distintos são
verificáveis. Fonte: \src{tests/castle/run\_castle\_evaluation.sh}
(\code{CWE\_MODE}) e \src{tests/juliet/run\_juliet\_evaluation.sh:108--126}.
\enquote{---} significa que o CWE não está naquele benchmark, não que seja
insuportado.}\label{tab:cwe}\\
\toprule
\textbf{CWE} & \textbf{Descrição} & \textbf{Modo} & \textbf{CASTLE} & \textbf{Juliet} \\
\midrule
\endfirsthead
\multicolumn{5}{@{}l}{\small\emph{(continuação da Tabela \thetable)}}\\
\toprule
\textbf{CWE} & \textbf{Descrição} & \textbf{Modo} & \textbf{CASTLE} & \textbf{Juliet} \\
\midrule
\endhead
\midrule
\multicolumn{5}{r@{}}{\small\emph{continua na próxima página}}\\
\endfoot
\bottomrule
\endlastfoot
\multicolumn{5}{@{}l}{\emph{Segurança de memória --- \flag{memtrack}}}\\
\midrule
121 & Stack-based buffer overflow          & \flag{memtrack} & --- & \yes \\
122 & Heap-based buffer overflow           & \flag{memtrack} & --- & \yes \\
125 & Out-of-bounds read                   & \flag{memtrack} & \yes & --- \\
415 & Double free                          & \flag{memtrack} & \yes & \yes \\
416 & Use after free                       & \flag{memtrack} & \yes & \yes \\
476 & NULL pointer dereference             & \flag{memtrack} & \yes & \yes \\
761 & Free of pointer not at start of buffer & \flag{memtrack} & \yes & \yes \\
787 & Out-of-bounds write                  & \flag{memtrack} & \yes & --- \\
822 & Untrusted pointer dereference        & \flag{memtrack} & \yes & --- \\
843 & Type confusion (\emph{access of resource with incompatible type}) & \flag{memtrack} & \yes & --- \\
\midrule
\multicolumn{5}{@{}l}{\emph{Aritmética --- \flag{check-overflow}}}\\
\midrule
190 & Integer overflow / wraparound        & \flag{check-overflow} & \yes & \yes \\
191 & Integer underflow                    & \flag{check-overflow} & --- & \yes \\
369 & Divide by zero                       & \flag{check-overflow} & \yes & \yes \\
\midrule
\multicolumn{5}{@{}l}{\emph{Ciclo de vida de memória --- \flag{memcleanup-property}}}\\
\midrule
401 & Missing release of memory (\emph{memory leak}) & \flag{memcleanup-property} & \yes & \yes \\
\midrule
\multicolumn{5}{@{}l}{\emph{Assertivas --- \flag{check-asserts}}}\\
\midrule
617 & Reachable assertion                  & \flag{check-asserts} & \yes & --- \\
\end{longtable}

\subsection{CWEs explicitamente fora de escopo}

O runner do CASTLE declara treze CWEs como fora de escopo. Isso é uma decisão
deliberada e testada, não uma omissão: \src{test\_cwe\_mode\_mapping.sh}
\emph{exige} que todo CWE do benchmark esteja ou mapeado ou declarado fora de
escopo, justamente para que nada caia num modo degenerado em silêncio.

\begin{table}[H]
\centering
\small
\begin{tabularx}{\textwidth}{@{}L{4.5cm} Y@{}}
\toprule
\textbf{CWEs} & \textbf{Motivo} \\
\midrule
628, 674, 770, 835 &
Terminação e consumo de recurso (recursão descontrolada, alocação sem limite,
laço infinito). O Map2Check não implementa a propriedade de terminação. \\
\addlinespace
22, 78, 89, 134 &
Injeção e uso indevido de dados externos (\emph{path traversal}, injeção de
comando/SQL, \emph{format string}). Exigem modelo de \emph{taint}, ausente. \\
\addlinespace
253, 327, 362, 522, 798 &
Verificação incorreta de retorno, criptografia fraca, condição de corrida,
credenciais desprotegidas/embutidas. Nenhuma é expressável como uma das cinco
propriedades implementadas. \\
\bottomrule
\end{tabularx}
\caption{Fonte: \src{run\_castle\_evaluation.sh:49}
(\code{OUT\_OF\_SCOPE\_CWES}).}
\end{table}

\subsection{Cobertura no modo WebAssembly}

Com \flag{wasm} a entrada é um binário \code{.wasm} elevado a C via
\code{wasm2c} (WABT $\geq$ 1.0.41) e daí a LLVM-IR. A análise é
\emph{agnóstica à linguagem-fonte}: instrumenta a \emph{linear memory} e as
chamadas \code{wasm\_rt\_*}, comuns a todo módulo elevado. O modo é forçado
para \code{MEMTRACK\_MODE}.

\begin{table}[H]
\centering
\small
\begin{tabularx}{\textwidth}{@{}Y L{3.6cm} C{2.2cm}@{}}
\toprule
\textbf{Propriedade} & \textbf{CWEs de exemplo} & \textbf{Estado} \\
\midrule
Acesso fora dos limites da \emph{linear memory}\newline
{\footnotesize (reportado como \code{FALSE-OVERFLOW}, não \code{FALSE-DEREF})}
& 121, 122, 124, 126, 127 & \yes \\
Overflow aritmético & 190 & \yes \\
Alcançabilidade de função-alvo/erro & --- & \yes \\
\emph{Leak} / \emph{use-after-free} / \emph{double free} & 401, 416, 415 & \no \\
\bottomrule
\end{tabularx}
\caption{Fonte: \src{README.md:92--99}. A última linha é trabalho futuro por uma
razão estrutural: no código elevado, o alocador da linguagem-fonte é compilado
\emph{dentro} da linear memory, então as chamadas \code{malloc}/\code{free} não
existem no IR. Rastrear memória ali exigiria uma convenção de exportação ou
heurística de detecção de alocador.}
\end{table}

%================================================================
\section{Eixos ortogonais}

Não são propriedades, mas multiplicam a matriz de execução e precisam ser
fixados em qualquer comparação para que ela seja justa.

\begin{table}[H]
\centering
\small
\begin{tabularx}{\textwidth}{@{}L{4.4cm} Y@{}}
\toprule
\textbf{Eixo} & \textbf{Valores} \\
\midrule
Gerador de \emph{nondet} & \code{fuzzer} (LibFuzzer), \code{symex} (KLEE), nenhum \\
Solucionador SMT & \code{z3} (padrão) e \code{stp} passam direto ao KLEE;
\code{btor} e \code{yices2} são roteados por \code{metasmt} e \textbf{não
funcionam} na imagem de referência (ver Seção~\ref{sec:limitacoes}, item~9) \\
Entrada & C / \code{.i} / \code{.bc}; ou \code{.wasm} via \flag{wasm} \\
Função de entrada & \flag{entry-function} (padrão \code{main}) \\
Modificadores & \flag{slice}, \flag{seed-exchange}, \flag{add-invariants}, \flag{btree} \\
Saídas & \flag{generate-witness}, \flag{generate-test-suite}, \flag{generate-testcase}, \flag{print-counter} \\
\bottomrule
\end{tabularx}
\caption{Fonte: \src{map2check.cpp:719--766}.}
\end{table}

\paragraph{Funções \emph{nondet} reconhecidas (16).}
\code{\_\_VERIFIER\_nondet\_} seguido de:
\code{int}, \code{uint}, \code{unsigned}, \code{char}, \code{uchar},
\code{pchar}, \code{short}, \code{ushort}, \code{long}, \code{ulong},
\code{bool}, \code{pointer}, \code{size\_t}, \code{loff\_t},
\code{sector\_t}, \code{double}
(\src{pass/NonDetPass.cpp:155--171}). Além disso,
\code{\_\_VERIFIER\_assume} e \code{assume\_abort\_if\_not} são reescritos como
poda de caminho em vez de \code{abort()}
(\src{NonDetPass.cpp:33--99}) --- sem isso, uma suposição falhando deixava um
\code{abort.err} indistinguível de uma violação real.

\paragraph{Códigos de saída.}
\code{0} sucesso; \code{1} erro de linha de comando; \code{3} a opção é válida
mas esta compilação não tem a capacidade (é o caso de \flag{add-invariants}
sem \code{-DENABLE\_CLAM=ON}); \code{4} o veredicto divergiu de
\flag{expected-result}. O código \code{4} substituiu um \code{abort()}, que
fazia o BenchExec ler \enquote{resposta errada} como \enquote{ferramenta
quebrada} (\src{map2check.cpp:47--60}).

%================================================================
\section{Limitações conhecidas}
\label{sec:limitacoes}

Registradas aqui porque uma comparação que as ignore mede a coisa errada.

\begin{enumerate}[leftmargin=*,itemsep=3pt]
  \item \textbf{Overflow sem sinal não é instrumentado.} As variantes
        \code{uint} de \code{Add}/\code{Sub}/\code{Mul} estão no código, porém
        comentadas. Isso está alinhado com a definição do SV-COMP, mas diverge
        do que uma ferramenta genérica de análise poderia reportar.
  \item \textbf{Ponto flutuante é \emph{no-op} deliberado.}
        \code{FAdd}/\code{FSub}/\code{FMul}/\code{FDiv}/\code{FRem} não geram
        veredicto.
  \item \textbf{Deslocamento à esquerda não é checado.} O \code{case} de
        \code{Shl} em \src{OverflowPass.cpp:395} tem corpo vazio, como
        \code{LShr}, \code{AShr}, \code{And}, \code{Or} e \code{Xor}. Um
        \code{x << n} que perde bits significativos passa como \emph{TRUE}.
  \item \textbf{Divisão por zero não tem flag própria} --- só é alcançada via
        \flag{check-overflow}.
  \item \textbf{Três veredictos geram uma \emph{witness} de corretude apesar de
        serem violações.} O \code{switch} de
        \src{witness/graph.cpp:208--237} tem \code{case} para
        \code{FALSE-FREE}, \code{FALSE-DEREF}, \code{FALSE-MEMTRACK},
        \code{TARGET-REACHED} e \code{FALSE-OVERFLOW}. \code{ASSERT},
        \code{FALSE-MEMCLEANUP} e \code{FALSE-DIVBYZERO} caem no
        \code{default}, que constrói um \code{CorrectnessWitnessGraph} --- e,
        sem \code{spectTrue}, o rotula com a especificação \code{MEMSAFETY}.
        Afeta apenas execuções com \flag{generate-witness}; o veredicto
        impresso e o código de saída continuam corretos.
  \item \textbf{Alcançabilidade com alvo \code{main} é degenerada.} O
        \code{TargetPass} instrumenta \emph{sítios de chamada} cuja função
        chamada tem o nome dado; um programa nunca chama o próprio \code{main}.
        Na baseline v6, 98 de 217 execuções passaram por esse modo e produziram
        zero \code{FALSE} em 59 programas sabidamente vulneráveis. O teste
        \src{test\_cwe\_mode\_mapping.sh} hoje proíbe esse mapeamento.
  \item \textbf{Invariantes (Clam) são \emph{opt-in}.} Um invariante não-sólido
        produz um \code{TRUE} errado em vez de um erro, o que é a falha mais
        cara possível --- por isso exige \code{-DENABLE\_CLAM=ON} \emph{e} o
        driver instalado.
  \item \textbf{\emph{Slicing}, \emph{seed-exchange} e WASM rodam nos defaults}
        da biblioteca; estão integrados, não ajustados. Qualquer número medido
        sobre eles reflete a configuração padrão, não um limite da técnica.
  \item \textbf{Dois dos quatro solucionadores não existem no build.} A CLI
        aceita \code{z3}, \code{stp}, \code{btor} e \code{yices2}
        (\src{map2check.cpp:804}), mas \code{btor} e \code{yices2} são
        entregues ao KLEE como \code{-{}-solver-backend=metasmt}
        (\src{caller.cpp:748--758}) e o \code{Dockerfile.dev} compila o KLEE
        apenas com \code{ENABLE\_SOLVER\_Z3} e \code{ENABLE\_SOLVER\_STP} ---
        sem metaSMT, sem Boolector, sem Yices. Passar um deles é aceito na
        linha de comando e falha na execução.
  \item \textbf{O veredicto de overflow é gravado como \code{OVERFLOW}, não
        \code{FALSE-OVERFLOW}.} \src{PropertyGenerator.c:97} escreve a string
        \code{OVERFLOW} em \code{map2check\_property}, e é isso que
        \src{tools.cpp:228} lê de volta. \code{FALSE-OVERFLOW} é o nome do
        valor do \emph{enum}, não o texto no arquivo: um harness que procure
        \code{FALSE-OVERFLOW} não casa com nada.
  \item \textbf{\code{UNKNOWN} nunca é gravado, e o arquivo alternativo tem o
        nome errado.} O \code{switch} de \src{PropertyGenerator.c} tem
        \code{case UNKNOWN:}, mas as duas únicas chamadas
        \code{write\_property(UNKNOWN, ...)} estão comentadas
        (\src{AnalysisModeMemtrack.c:308},
        \src{AnalysisModeMemcleanup.c:298}) --- é código morto, do mesmo tipo
        do \code{case} de \code{Shl}. O runtime usa
        \code{write\_property\_unknown()}, que escreve num arquivo à parte
        chamado \code{map2check\_property\_\_unknown} (dois sublinhados,
        \src{PropertyGenerator.c:16}); o frontend abre
        \code{map2check\_property\_unknown} (um sublinhado,
        \src{tools.cpp:114}). Os dois nomes não coincidem, então esse arquivo
        também nunca é lido. O tool ainda \emph{reporta}
        \code{VERIFICATION UNKNOWN}, mas por outras rotas --- o inicializador
        do campo (\src{tools.hpp:136}), o \code{map2check\_property} ausente ou
        ilegível (\src{tools.cpp:102,252}) e o próprio frontend em caso de
        estouro de tempo (\src{map2check.cpp:581,585}). O que não existe é o
        caminho pelo qual o \emph{runtime} declararia indecisão.
  \item \textbf{KLEE: \flag{max-time} exige segundos inteiros.} Um valor
        fracionário encerra a execução em silêncio.
\end{enumerate}

%================================================================
\section{Resumo quantitativo}

\begin{table}[H]
\centering
\small
\begin{tabularx}{\textwidth}{@{}Y c@{}}
\toprule
\textbf{Dimensão} & \textbf{Total} \\
\midrule
Modos de verificação implementados & 6 \\
\quad dos quais checam uma propriedade & 5 \\
Valores do \emph{enum} \code{ViolatedProperty} & 10 \\
\quad realmente gravados em \code{map2check\_property} & 9 \\
\quad que são violações concretas & 8 \\
\quad que geram \emph{witness} de violação & 5 \\
Tipos de especificação de \emph{witness} (\code{SpecificationType}) & 6 \\
Instruções LLVM instrumentadas para overflow & 5 \\
Funções \emph{nondet} reconhecidas & 16 \\
CWEs verificáveis (união CASTLE $\cup$ Juliet) & 15 \\
CWEs declarados fora de escopo & 13 \\
Solucionadores SMT aceitos pela CLI & 4 \\
\quad utilizáveis na imagem de referência (Z3, STP) & 2 \\
Geradores de entrada selecionáveis (\flag{nondet-generator}) & 2 \\
\bottomrule
\end{tabularx}
\end{table}

\end{document}
